\documentclass[11pt,a4paper]{article}

% --- Robust pdfLaTeX preamble (Times-like, no fontspec dependency) ---
\usepackage[T1]{fontenc}
\usepackage[utf8]{inputenc}
\usepackage[british]{babel}
\usepackage[round,authoryear]{natbib}
\usepackage{newtxtext,newtxmath}
\usepackage[margin=25mm]{geometry}
\usepackage{microtype}
\usepackage{booktabs}
\usepackage{array}
\usepackage{enumitem}
\usepackage{titlesec}
\usepackage{xcolor}
\usepackage[hidelinks]{hyperref}

\frenchspacing

\setlist[itemize]{leftmargin=1.4em, topsep=3pt, itemsep=2pt, parsep=0pt}
\setlist[enumerate]{leftmargin=1.8em, topsep=3pt, itemsep=2pt, parsep=0pt}
\setlength{\parskip}{4pt}
\setlength{\parindent}{0pt}
\renewcommand{\arraystretch}{1.15}

\definecolor{rule}{gray}{0.45}
\titleformat{\section}{\large\bfseries\sffamily}{\thesection.}{0.6em}{}
\titleformat{\subsection}{\normalsize\bfseries\sffamily}{\thesubsection}{0.6em}{}
\titlespacing*{\section}{0pt}{14pt}{4pt}
\titlespacing*{\subsection}{0pt}{8pt}{2pt}

\newcommand{\fname}[1]{\texttt{\small #1}}

\begin{document}

\begin{center}
{\LARGE\bfseries\sffamily Assessment of the proposed \textit{Metacommunity\\ Ecology} chapter for BCB743}\\[6pt]
{\large\sffamily and its relation to a planned chapter on niche and neutral theory}\\[10pt]
{\normalsize A.\,J.\ Smit \quad$\cdot$\quad Module coordinator, BCB743 Quantitative Ecology}\\[2pt]
{\small 22 June 2026}
\end{center}

\vspace{2pt}
\hrule height 0.8pt
\vspace{6pt}

{\small\itshape Scope. This memo assesses the chapter-contribution presentation
\fname{metacommunity\_ecology.pptx} (Task A2, presentation component) by A.\ Waglay,
B.\ Effendi, and C.\ Mc Intyre. It evaluates the proposal against the existing BCB743
material and against a separate, planned chapter on niche and neutral theory, and it
records design recommendations for both chapters.}

\vspace{6pt}
\hrule height 0.4pt
\vspace{8pt}

\section{Summary and recommendation}

The presentation identifies a real and specific gap, and it proposes a sound device for
closing it. BCB743 teaches a multivariate toolbox (ordination, dissimilarity methods,
constrained ordination, regression) with considerable rigour, but the ecological theory
that gives those tools their purpose is currently scattered across asides and worked
examples rather than developed in its own right. The metacommunity framework is a
defensible way to organise that theory around the methods already on the syllabus.

Three judgements follow.

\begin{itemize}
  \item \textbf{Accept the proposal in principle, with revision.} The diagnosis on
  slide~2, namely that the existing tools are not unified under a single framework that
  explains why communities differ across space, is accurate and aligns with the module's
  own stated outcomes.
  \item \textbf{Design the two chapters together, not separately.} The proposed
  metacommunity chapter and the planned niche/neutral chapter describe the same
  underlying processes at two levels of organisation. Treated as independent units they
  will overlap and, worse, may present niche and neutral ideas as competing frameworks
  when the current synthesis treats them as endpoints of one process continuum.
  \item \textbf{Sequence theory first, spatial operationalisation second.} A short
  niche/neutral chapter should establish the conceptual foundation early; the
  metacommunity chapter should follow later as the spatial application that connects to
  variation partitioning, after db-RDA has been taught.
\end{itemize}

\section{What the presentation proposes}

The deck is a nine-slide conceptual overview. It states the gap (slide~2), defines a
metacommunity as a set of local communities connected by dispersal of multiple
interacting species (slide~3), lists the four processes that generate community
differences (slide~4: environmental filtering, dispersal limitation, species
interactions, stochasticity), and presents the four metacommunity paradigms of
\citet{leibold2004} (slide~5: species sorting, mass effects, neutral theory, patch
dynamics). It then maps each paradigm onto methods already in the module (slide~6) and
sketches a worked example on the South African seaweed data (slide~7: Bray--Curtis,
nMDS, db-RDA, variation partitioning). The closing argument (slide~8) is that the
chapter would integrate community and spatial ecology, connect niche and neutral theory,
and reuse existing module datasets and methods. The deck carries no equations, no R code,
and no worked output; it is a proposal sketch rather than a chapter draft.

\section{Is the gap real?}

Yes, and it is more precise than the deck states. The module's published outcomes already
commit to this material. The index lists, under the core theoretical framework,
``ecological hypotheses about how species assemble across space and time, including
niche-based mechanisms, neutral processes, and historical contingencies'', and one
outcome requires students to ``know and understand the current hypotheses that explain
species assembly processes in space and time (e.g.\ neutral and niche mechanisms)''. No
current chapter develops that theory. The word \emph{metacommunity} appears nowhere in
the syllabus.

Where the theory does surface, it does so as an aside to a method. The db-RDA chapter
(\fname{constrained\_ordination.qmd}) frames the seaweed result as a niche-difference
mechanism operating under Beijerinck's ``everything is everywhere, but the environment
selects'', which is a species-sorting interpretation in all but name. The \fname{spesim}
chapter exposes the generating processes directly, namely a species-abundance
distribution (including a neutral zero-sum multinomial sampler), environmental filtering
as a Gaussian response on a gradient, and a dispersal kernel. The raw ingredients of
assembly theory are therefore present in the course, but they are not assembled. The
proposal is a request for that synthesis, and the request is justified.

\section{The conceptual relationship: niche, neutral, and the metacommunity paradigms}

This is the decisive point for designing both chapters, and it constrains how they should
be written.

Niche theory and neutral theory are not opposed accounts of the same phenomenon; they are
the endpoints of a continuum of process. \citet{vellend2010,vellend2016} reduces
community dynamics to four high-level processes, namely \emph{selection} (deterministic
fitness differences among species, i.e.\ niche-based sorting), \emph{drift} (stochastic
change in abundances), \emph{dispersal} (movement among localities), and
\emph{speciation}. Niche theory is the limit in which selection dominates; neutral theory
\citep{hubbell2001} is the limit in which species are assumed ecologically equivalent, so
that drift and dispersal alone govern composition. Most real assemblages sit between
these limits \citep{gravel2006,chase2011}.

The four metacommunity paradigms are positions in that same process space, not a separate
theory. Species sorting is selection with moderate dispersal. Mass effects add high
dispersal, so that source--sink flux overrides local selection. Patch dynamics turn on
colonisation--extinction trade-offs among competitors. The neutral paradigm is drift plus
dispersal under species equivalence. \citet{leibold2004} introduced the four-way typology;
\citet{leibold2018} and the empirical review of \citet{logue2011} reframe the same ground
in terms of underlying processes rather than discrete archetypes.

The implication for the curriculum is direct. The niche/neutral chapter and the
metacommunity chapter are the same content at two altitudes. The first should establish
the processes; the second should show how those processes play out across space and how
they are tested. If they are written independently, they will duplicate the neutral-theory
exposition and risk teaching the four paradigms as rival hypotheses rather than as a
cross-classification of selection, drift, and dispersal.

\section{Fit with the existing methods curriculum}

The method mapping on slide~6 is largely sound. PCA recovers environmental gradients; nMDS
and PCoA summarise compositional similarity and $\beta$-diversity; cluster analysis groups
sites; regression tests environmental drivers. The load-bearing link, however, is
variation partitioning, and here the proposal connects to the course more deeply than the
students appear to realise.

The standard quantitative test for niche versus neutral signal in a metacommunity is the
partition of community variation into an environmental fraction~$[E]$, a spatial
fraction~$[S]$, their shared component, and a residual \citep{borcard1992,
peresneto2006, cottenie2005}. That exact analysis already exists in the module. The
self-study appendix \fname{two\_oceans\_appendices.qmd} (Chapter~13b) builds spatial
eigenfunctions with \fname{PCNM()}, then calls \fname{vegan::varpart()} to split seaweed
$\beta$-diversity into thermal and spatial fractions, citing \citet{peresneto2006} and
\citet{peresneto2010}. The lectured db-RDA chapter, by contrast, fits the environmental
model only and explicitly sets the spatial coordinates aside.

So the method the proposal needs is taught, but as an unmotivated technical appendix. The
metacommunity chapter would give that orphaned analysis its ecological reason for
existing, namely separating the environmental (species-sorting) signal from the spatial
(dispersal-related) signal. That is a strong argument for adding the chapter, and a clean
integration point: the chapter does not introduce a new method so much as supply the
theory that makes an existing one interpretable.

One terminological hazard needs managing. The course already uses the word
\emph{partition} in a different sense. Chapter~13a-revised (\fname{constrained\_ordination\_v2.qmd})
partitions $\beta$-diversity \emph{dissimilarity} into turnover and nestedness components
\citep{baselga2010}, and then into replacement and richness-difference components
\citep{legendre2014}. That is a decomposition of a dissimilarity coefficient, not the
environment--space decomposition of community variation. The two operations share a word
and nothing else. The new chapter must keep the two senses distinct or it will confuse
students who have just met the first.

\section{Critical caveats}

A chapter built on this framework must assess its own tools, not advertise them. Three
caveats should be written into the chapter rather than left for the lecturer to add.

\subsection*{(a) The four-paradigm typology is dated}
The discrete four-way scheme of \citet{leibold2004} has largely given way to the
process-based view \citep{vellend2010,logue2011,leibold2018}. The paradigms remain a
useful teaching scaffold, but presenting them as the current state of the field would
misrepresent it. Teach the typology as history, then move to selection--drift--dispersal
as the operative framework.

\subsection*{(b) Variation partitioning is a weak test of process}
The $[E]$ versus $[S]$ partition is suggestive, not diagnostic, and the chapter must say
so. Environmental variables are themselves spatially structured, so the shared fraction is
typically large and hard to interpret. The pure spatial fraction~$[S]$ conflates dispersal
processes with spatially structured but unmeasured environment, since spatial
eigenfunctions absorb autocorrelation from any source \citep{gilbert2010, smith2010,
tuomisto2006, peresneto2010}. Reporting~$[S]$ as ``the neutral component'' would be
indefensible. The honest claim is that a partition is consistent with a relative balance
of processes, under stated assumptions, and requires further testing.

\subsection*{(c) The seaweed dataset suits one paradigm, not all four}
The South African seaweed data describe an essentially one-dimensional thermal gradient
sampled by contiguous coastal sections. That design is well suited to demonstrating
species sorting, which is exactly the conclusion the existing db-RDA chapter reaches under
Beijerinck's law. It is poorly suited to demonstrating mass effects or patch dynamics,
which require patchy, dispersal-structured landscapes with replicated patches. The worked
example should therefore frame the seaweeds as a species-sorting case study and extend the
existing interpretation, not promise a demonstration of all four paradigms it cannot
deliver.

\section{Recommended sequencing and division of labour}

The two chapters should form a conceptual pair, placed at either side of the constrained
ordination block.

\begin{center}
\begin{tabular}{@{}>{\raggedright\arraybackslash}p{0.30\linewidth} >{\raggedright\arraybackslash}p{0.62\linewidth}@{}}
\toprule
\textbf{Chapter} & \textbf{Content and placement}\\
\midrule
Niche \& neutral theory \emph{(new, conceptual)} &
Selection, drift, dispersal, and speciation as the processes of community assembly
\citep{vellend2010}; the niche concept and modern coexistence theory
\citep{chase2003, hillerislambers2012}; Hubbell's neutral theory and its testable
predictions for the species-abundance distribution, dispersal limitation, and distance
decay of similarity \citep{hubbell2001}; the niche--neutral continuum. \emph{Place early},
after the introduction to ordination and $\beta$-diversity, to motivate why turnover is
measured at all.\\
\addlinespace
Metacommunity ecology \emph{(the proposal)} &
The four paradigms as positions in selection--drift--dispersal space; variation
partitioning ($[E]$ vs $[S]$) as the quantitative test and its limitations; a worked
seaweed example that reuses the Chapter~13b PCNM and \fname{varpart()} machinery.
\emph{Place late}, after db-RDA and variation partitioning are in hand.\\
\bottomrule
\end{tabular}
\end{center}

This bookend arrangement closes the standing gap between the module's stated outcomes and
its delivered content, and it gives the self-study variation-partitioning appendix a
lectured home.

\section{An integration the proposal misses: \fname{spesim}}

The \fname{spesim} package is the controlled setting in which the caveats above become a
practical exercise. It already simulates the three processes the chapters turn on, namely a
species-abundance distribution (including a neutral zero-sum multinomial), environmental
filtering as a Gaussian response on a gradient, and dispersal through a kernel. Students
can generate communities under a \emph{known} balance of species sorting and neutral
drift, run the $[E]$/$[S]$ partition, and check whether it recovers the truth they
specified. That exercise demonstrates caveat~(b) directly rather than asserting it, and it
triangulates with Task~A3, the \fname{spesim} package review. The proposal should adopt it.

\section{The presentation as an assessed artefact (Task A2)}

\textbf{Strengths.} The gap is correctly diagnosed and tied to the module's own framing.
The method mapping is accurate. The chosen dataset reuses material already in the course,
which is the right instinct for a teaching chapter. The structure moves cleanly from
problem to framework to application.

\textbf{Weaknesses relative to a deliverable chapter.} The deck is conceptual only, with no
equations, no R, and no output, so the distance to a worked BCB743 chapter is large. The
four paradigms are presented without the modern synthesis or the inferential caveats.
``Variation partitioning'' is invoked on slide~7 without being shown and without being
distinguished from the $\beta$-diversity dissimilarity partition the students will already
have met. The \fname{spesim} opportunity is absent.

\textbf{Verdict.} A sound and well-judged proposal. The framework is appropriate, the fit
with the module is genuine, and the gap it targets is real. The work that remains is
substantial and lies mainly in the worked quantitative example and in the critical framing
of what variation partitioning can and cannot establish.

\vspace{8pt}
\hrule height 0.4pt
\vspace{6pt}

\renewcommand{\refname}{\normalsize\sffamily\bfseries References}
\begin{thebibliography}{99}
\small
\setlength{\itemsep}{1pt}

\bibitem[Baselga(2010)]{baselga2010} Baselga, A. (2010) Partitioning the turnover and
nestedness components of beta diversity. \textit{Global Ecology and Biogeography}, 19,
134--143.

\bibitem[Borcard et al.(1992)]{borcard1992} Borcard, D., Legendre, P. \& Drapeau, P.
(1992) Partialling out the spatial component of ecological variation. \textit{Ecology},
73, 1045--1055.

\bibitem[Chase \& Leibold(2003)]{chase2003} Chase, J.M. \& Leibold, M.A. (2003)
\textit{Ecological Niches: Linking Classical and Contemporary Approaches}. University of
Chicago Press, Chicago.

\bibitem[Chase \& Myers(2011)]{chase2011} Chase, J.M. \& Myers, J.A. (2011) Disentangling
the importance of ecological niches from stochastic processes across scales.
\textit{Philosophical Transactions of the Royal Society B}, 366, 2351--2363.

\bibitem[Cottenie(2005)]{cottenie2005} Cottenie, K. (2005) Integrating environmental and
spatial processes in ecological community dynamics. \textit{Ecology Letters}, 8,
1175--1182.

\bibitem[Gilbert \& Bennett(2010)]{gilbert2010} Gilbert, B. \& Bennett, J.R. (2010)
Partitioning variation in ecological communities: do the numbers add up? \textit{Journal
of Applied Ecology}, 47, 1071--1082.

\bibitem[Gravel et al.(2006)]{gravel2006} Gravel, D., Canham, C.D., Beaudet, M. \&
Messier, C. (2006) Reconciling niche and neutrality: the continuum hypothesis.
\textit{Ecology Letters}, 9, 399--409.

\bibitem[HilleRisLambers et al.(2012)]{hillerislambers2012} HilleRisLambers, J., Adler,
P.B., Harpole, W.S., Levine, J.M. \& Mayfield, M.M. (2012) Rethinking community assembly
through the lens of coexistence theory. \textit{Annual Review of Ecology, Evolution, and
Systematics}, 43, 227--248.

\bibitem[Hubbell(2001)]{hubbell2001} Hubbell, S.P. (2001) \textit{The Unified Neutral
Theory of Biodiversity and Biogeography}. Princeton University Press, Princeton.

\bibitem[Legendre(2014)]{legendre2014} Legendre, P. (2014) Interpreting the replacement
and richness difference components of beta diversity. \textit{Global Ecology and
Biogeography}, 23, 1324--1334.

\bibitem[Leibold et al.(2004)]{leibold2004} Leibold, M.A., Holyoak, M., Mouquet, N.,
Amarasekare, P., Chase, J.M., Hoopes, M.F., Holt, R.D., Shurin, J.B., Law, R., Tilman, D.,
Loreau, M. \& Gonzalez, A. (2004) The metacommunity concept: a framework for multi-scale
community ecology. \textit{Ecology Letters}, 7, 601--613.

\bibitem[Leibold \& Chase(2018)]{leibold2018} Leibold, M.A. \& Chase, J.M. (2018)
\textit{Metacommunity Ecology}. Princeton University Press, Princeton.

\bibitem[Logue et al.(2011)]{logue2011} Logue, J.B., Mouquet, N., Peter, H. \&
Hillebrand, H. (2011) Empirical approaches to metacommunities: a review and comparison
with theory. \textit{Trends in Ecology \& Evolution}, 26, 482--491.

\bibitem[Peres-Neto et al.(2006)]{peresneto2006} Peres-Neto, P.R., Legendre, P., Dray, S.
\& Borcard, D. (2006) Variation partitioning of species data matrices: estimation and
comparison of fractions. \textit{Ecology}, 87, 2614--2625.

\bibitem[Peres-Neto \& Legendre(2010)]{peresneto2010} Peres-Neto, P.R. \& Legendre, P.
(2010) Estimating and controlling for spatial structure in the study of ecological
communities. \textit{Global Ecology and Biogeography}, 19, 174--184.

\bibitem[Smith \& Lundholm(2010)]{smith2010} Smith, T.W. \& Lundholm, J.T. (2010)
Variation partitioning as a tool to distinguish between niche and neutral processes.
\textit{Ecography}, 33, 648--655.

\bibitem[Tuomisto \& Ruokolainen(2006)]{tuomisto2006} Tuomisto, H. \& Ruokolainen, K.
(2006) Analyzing or explaining beta diversity? Understanding the targets of different
methods of analysis. \textit{Ecology}, 87, 2697--2708.

\bibitem[Vellend(2010)]{vellend2010} Vellend, M. (2010) Conceptual synthesis in community
ecology. \textit{The Quarterly Review of Biology}, 85, 183--206.

\bibitem[Vellend(2016)]{vellend2016} Vellend, M. (2016) \textit{The Theory of Ecological
Communities}. Princeton University Press, Princeton.

\end{thebibliography}

\end{document}
